% Part: normal-modal-logic
% Chapter: sequent-calculus
% Section: more-rules
%READERNOTE{SOL6-A185: मोडल क्रमवर्ती प्रणाल्यांच्या तक्त्यात मूळ नियमांची नावे primitive Box/Diamond निवडीनुसार केली; दोन्ही primitive असताना दोन्ही नियम दाखवले. Diamond-only कट-सिद्धतेतील डाव्या सूत्रांच्या अदलाबदलीचे चुकीचे उजवे नियम-नाव डावे केले. कटविना असंपूर्णतेचा दावा या विभागातील विशिष्ट कलनापुरता स्पष्ट केला.}

\documentclass[../../../include/open-logic-section]{subfiles}

\begin{document}

\olfileid{nml}{seq}{mru}

\olsection{इतर प्राप्यता-संबंधांसाठी नियम}

विशिष्ट प्राप्यता-संबंधांनी निर्धारित होणाऱ्या
तर्कशास्त्रांचा विचार करण्यासाठी आपण
\olref{tab:more-rules} मधील अतिरिक्त नियम पाहू.

\iftag{prvBox}{\iftag{prvDiamond}{% [] and <> primitive
\begin{table}
  \begin{center}
    \def\arraystretch{3}
    \begin{tabular}{cc}
    \hline
      \Axiom$!A, \Gamma \fCenter \Delta$
      \RightLabel{T$\Box$}
      \UnaryInf$\Box!A, \Gamma \fCenter \Delta$
      \DisplayProof
    &
      \Axiom$\Gamma \fCenter \Delta, !A$
      \RightLabel{T$\Diamond$}
      \UnaryInf$\Gamma \fCenter \Delta, \Diamond!A$
      \DisplayProof
    \\[1ex]
%    \hline
    \multicolumn{2}{c}{
      \Axiom$\Gamma \fCenter \Delta$
      \RightLabel{D}
      \UnaryInf$\Box\Gamma \fCenter \Diamond\Delta$
      \DisplayProof}
    \\[1ex]
%    \hline
      \Axiom$\Gamma, \Diamond\Pi \fCenter \Box\Delta, \Lambda, !A$
      \RightLabel{B$\Box$}
      \UnaryInf$\Box\Gamma, \Pi \fCenter \Delta, \Diamond\Lambda, \Box!A$
      \DisplayProof
      &
      \Axiom$!A, \Diamond\Gamma, \Pi \fCenter \Box\Lambda, \Delta$
      \RightLabel{B$\Diamond$}
      \UnaryInf$\Diamond!A, \Gamma, \Box\Pi \fCenter \Lambda, \Diamond\Delta$
      \DisplayProof
    \\[1ex]
%    \hline
      \Axiom$\Box\Gamma \fCenter \Diamond\Delta, !A$
      \RightLabel{4$\Box$}
      \UnaryInf$\Box\Gamma \fCenter \Diamond\Delta, \Box!A$
      \DisplayProof
      &
      \Axiom$!A, \Box\Gamma \fCenter \Diamond\Delta$
      \RightLabel{4$\Diamond$}
      \UnaryInf$\Diamond!A, \Box\Gamma \fCenter \Diamond\Delta$
      \DisplayProof
    \\[1ex]
%    \hline
      \Axiom$\Box\Gamma, \Diamond\Pi \fCenter \Box\Delta,
      \Diamond\Lambda, !A$
      \RightLabel{5$\Box$}
      \UnaryInf$\Box\Gamma, \Diamond\Pi \fCenter \Box\Delta,
      \Diamond\Lambda, \Box!A$
      \DisplayProof
&
      \Axiom$!A, \Diamond\Gamma, \Box\Pi \fCenter \Diamond\Delta, \Box\Lambda$
      \RightLabel{5$\Diamond$}
      \UnaryInf$\Diamond!A,\Diamond\Gamma,\Box\Pi \fCenter \Diamond\Delta,\Box\Lambda$
      \DisplayProof
    \\[1ex]
    \hline
    \end{tabular}
  \end{center}
  \caption{अधिक मोडल नियम.}
  \ollabel{tab:more-rules}
\end{table}
}{ % only [] primitive
\begin{table}
  \begin{center}
    \def\arraystretch{3}
    \begin{tabular}{cc}
    \hline
      \Axiom$!A,\Gamma \fCenter \Delta$
      \RightLabel{T$\Box$}
      \UnaryInf$\Box!A, \Gamma \fCenter \Delta$
      \DisplayProof
    &
    %\hline
      \Axiom$\Gamma \fCenter$
      \RightLabel{D$\Box$}
      \UnaryInf$\Box\Gamma \fCenter$
      \DisplayProof
    \\[1ex]
    %\hline
      \Axiom$\Gamma \fCenter \Box\Delta, !A$
      \RightLabel{B$\Box$}
      \UnaryInf$\Box\Gamma \fCenter \Delta, \Box!A$
      \DisplayProof
    &
    %\hline
      \Axiom$\Box\Gamma \fCenter !A$
      \RightLabel{4$\Box$}
      \UnaryInf$\Box\Gamma \fCenter \Box!A$
      \DisplayProof
    \\[1ex]
    %\hline
    \multicolumn{2}{c}{\Axiom$\Box\Gamma \fCenter \Box\Delta, !A$
      \RightLabel{5$\Box$}
      \UnaryInf$\Box\Gamma \fCenter \Box\Delta, \Box!A$
      \DisplayProof}
    \\[1ex]
    \hline
    \end{tabular}
  \end{center}
  \caption{अधिक मोडल नियम.}
  \ollabel{tab:more-rules}
\end{table}
}}{% only <> primitive
\begin{table}
  \begin{center}
    \def\arraystretch{3}
    \begin{tabular}{cc}
    \hline
      \Axiom$\Gamma \fCenter \Delta, !A$
      \RightLabel{T$\Diamond$}
      \UnaryInf$\Gamma \fCenter \Delta, \Diamond!A$
      \DisplayProof
    &
    %\hline
      \Axiom$\fCenter \Delta$
      \RightLabel{D$\Diamond$}
      \UnaryInf$\fCenter \Diamond\Delta$
      \DisplayProof
    \\[1ex]
    %\hline
      \Axiom$!A, \Diamond\Gamma \fCenter \Delta$
      \RightLabel{B$\Diamond$}
      \UnaryInf$\Diamond!A, \Gamma \fCenter \Diamond\Delta$
      \DisplayProof
    &
    %\hline
      \Axiom$!A \fCenter \Diamond\Delta$
      \RightLabel{4$\Diamond$}
      \UnaryInf$\Diamond!A \fCenter \Diamond\Delta$
      \DisplayProof
    \\[1ex]
    %\hline
    \multicolumn{2}{c}{
      \Axiom$!A, \Diamond\Gamma \fCenter \Diamond\Delta$
      \RightLabel{5$\Diamond$}
      \UnaryInf$\Diamond!A,\Diamond\Gamma \fCenter \Diamond\Delta$
      \DisplayProof}
    \\[1ex]
    \hline
    \end{tabular}
  \end{center}
  \caption{अधिक मोडल नियम.}
  \ollabel{tab:more-rules}
\end{table}
}

हे नियम जोडल्यावर \olref{tab:logics-rules} मध्ये
दिलेल्या तर्कशास्त्रांसाठी निर्दोष आणि संपूर्ण
पद्धती मिळतात.

\begin{table}
  \begin{center}
    \begin{tabular}{lll}
      \hline
      तर्कशास्त्र & $R$ हा \dots आहे & नियम\\
      \hline
      $\Log{T} = \Log{KT}$ & परावर्ती & \iftag{prvBox}{$\Box$\iftag{prvDiamond}{, }{}}{}%
      \iftag{prvDiamond}{$\Diamond$}{},
      \iftag{prvBox}{T$\Box$}{}%
      \iftag{notprvBox,notprvDiamond}{}{, }%
      \iftag{prvDiamond}{T$\Diamond$}{}
      \\ \hline
      $\Log{D} = \Log{KD}$ & सर्वत्र उत्तरवर्ती & \iftag{prvBox}{$\Box$\iftag{prvDiamond}{, }{}}{}%
      \iftag{prvDiamond}{$\Diamond$}{},
      \iftag{prvBox}{%
        \iftag{prvDiamond}{D}{D$\Box$}}{D$\Diamond$}
      \\ \hline
      $\Log{K4}$ & संक्रमक & \iftag{prvBox}{$\Box$\iftag{prvDiamond}{, }{}}{}%
      \iftag{prvDiamond}{$\Diamond$}{},
      \iftag{prvBox}{4$\Box$}{}%
      \iftag{notprvBox,notprvDiamond}{}{, }%
      \iftag{prvDiamond}{4$\Diamond$}{}
      \\ \hline
      $\Log{B} = \Log{KTB}$ & परावर्ती, & \iftag{prvBox}{$\Box$\iftag{prvDiamond}{, }{}}{}%
      \iftag{prvDiamond}{$\Diamond$}{},
      \iftag{prvBox}{T$\Box$}{}%
      \iftag{notprvBox,notprvDiamond}{}{, }%
      \iftag{prvDiamond}{T$\Diamond$}{}\\
      & सममित &
      \iftag{prvBox}{B$\Box$}{}%
      \iftag{notprvBox,notprvDiamond}{}{, }%
      \iftag{prvDiamond}{B$\Diamond$}{}
      \\ \hline
      $\Log{S4} = \Log{KT4}$ & परावर्ती, & \iftag{prvBox}{$\Box$\iftag{prvDiamond}{, }{}}{}%
      \iftag{prvDiamond}{$\Diamond$}{},
      \iftag{prvBox}{T$\Box$}{}%
      \iftag{notprvBox,notprvDiamond}{}{, }%
      \iftag{prvDiamond}{T$\Diamond$}{}\\
      & संक्रमक &
      \iftag{prvBox}{4$\Box$}{}%
      \iftag{notprvBox,notprvDiamond}{}{, }%
      \iftag{prvDiamond}{4$\Diamond$}{}
      \\ \hline
      $\Log{S5} = \Log{KT5}$ & परावर्ती, & \iftag{prvBox}{$\Box$\iftag{prvDiamond}{, }{}}{}%
      \iftag{prvDiamond}{$\Diamond$}{},
      \iftag{prvBox}{T$\Box$}{}%
      \iftag{notprvBox,notprvDiamond}{}{, }%
      \iftag{prvDiamond}{T$\Diamond$}{}\\
      & संक्रमक, &
      \iftag{prvBox}{5$\Box$}{}%
      \iftag{notprvBox,notprvDiamond}{}{, }%
      \iftag{prvDiamond}{5$\Diamond$}{}\\
      &  युक्लिडी &
      \\ \hline
    \end{tabular}
  \end{center}
  \caption{विविध मोडल तर्कशास्त्रांसाठी क्रमवर्ती नियम.}
  \ollabel{tab:logics-rules}
\end{table}

\iftag{prvBox}{
\begin{ex}
  $\Log{K4} \Proves \Ax{4}$, म्हणजे
  $\Box!A \lif \Box\Box!A$, हे दाखवणारी क्रमवर्ती
  !!{derivation} पुढे दिली आहे.
  \begin{prooftree}
    \Axiom$\Box!A \fCenter \Box!A$
    \RightLabel{4$\Box$}
    \UnaryInf$\Box!A \fCenter \Box\Box!A$
    \RightLabel{\RightR{\lif}}
    \UnaryInf$\fCenter \Box!A \lif \Box\Box!A$
  \end{prooftree}
\end{ex}
}{\iftag{prvDiamond}{
\begin{ex}
  $\Log{K4} \Proves \Ax{4}$, म्हणजे
  $\Diamond\Diamond!A \lif \Diamond!A$, हे दाखवणारी
  क्रमवर्ती !!{derivation} पुढे दिली आहे.
  \begin{prooftree}
    \Axiom$\Diamond!A \fCenter \Diamond!A$
    \RightLabel{4$\Diamond$}
    \UnaryInf$\Diamond\Diamond!A \fCenter \Diamond!A$
    \RightLabel{\RightR{\lif}}
    \UnaryInf$\fCenter \Diamond\Diamond!A \lif \Diamond!A$
  \end{prooftree}
\end{ex}
}{}}

\iftag{prvBox}{\iftag{prvDiamond}{% <> and [] primitive
\begin{ex}
  $\Log{S5} \Proves \Ax{5}$, म्हणजे
  $\Diamond!A \lif \Box\Diamond!A$, हे दाखवणारी
  क्रमवर्ती !!{derivation} पुढे दिली आहे.
  \begin{prooftree}
    \Axiom$\Diamond!A \fCenter \Diamond!A$
    \RightLabel{5$\Box$}
    \UnaryInf$\Diamond!A \fCenter \Box\Diamond!A$
    \RightLabel{\RightR{\lif}}
    \UnaryInf$\fCenter \Diamond!A \lif \Box\Diamond!A$
  \end{prooftree}
\end{ex}
}{% only Box primitive
\begin{ex}
  $\Log{S5} \Proves \Ax{5}$, म्हणजे
  $\Diamond!A \lif \Box\Diamond!A$, हे दाखवणारी
  क्रमवर्ती !!{derivation} पुढे दिली आहे.
  \begin{prooftree}
    \Axiom$\Box\lnot!A \fCenter \Box\lnot!A$
    \RightLabel{\RightR{\lnot}}
    \UnaryInf$\fCenter \Box\lnot!A, \lnot\Box\lnot!A$
    \RightLabel{5$\Box$}
    \UnaryInf$ \fCenter \Box\lnot!A, \Box\lnot\Box\lnot!A$
    \RightLabel{\LeftR{\lnot}}
    \UnaryInf$\lnot\Box\lnot!A \fCenter \Box\lnot\Box\lnot!A$
    \RightLabel{\RightR{\lif}}
    \UnaryInf$\fCenter \lnot\Box\lnot!A \lif \Box\lnot\Box\lnot!A$
  \end{prooftree}
\end{ex}
}}{% only <> prmitive
\begin{ex}
  $\Log{S5} \Proves \Ax{5}$, म्हणजे
  $\Diamond!A \lif \Box\Diamond!A$, हे दाखवणारी
  क्रमवर्ती !!{derivation} पुढे दिली आहे.
  \begin{prooftree}
    \Axiom$\Diamond!A \fCenter\Diamond!A$
    \RightLabel{\LeftR{\lnot}}
    \UnaryInf$\lnot\Diamond !A, \Diamond!A \fCenter $
    \RightLabel{5$\Diamond$}
    \UnaryInf$\Diamond\lnot\Diamond!A, \Diamond!A \fCenter $
    \RightLabel{\RightR{\lnot}}
    \UnaryInf$\Diamond!A \fCenter \lnot\Diamond\lnot\Diamond!A$
    \RightLabel{\RightR{\lif}}
    \UnaryInf$\fCenter \Diamond!A \lif \lnot\Diamond\lnot\Diamond!A$
  \end{prooftree}
\end{ex}
}
\begin{ex}
या विभागात दिलेले \Log{S5} साठीचे क्रमवर्ती कलन \Cut{} नियमाशिवाय
संपूर्ण नाही; उदा., \iftag{prvBox}{$\Diamond\Box !A \lif !A$}{$!A \lif
\Box\Diamond !A$} हे~$\Log{S5}$ मध्ये वैध असलेले
सूत्र~\Cut शिवाय सिद्ध करता येत नाही. पुढे~\Cut
वापरणारी !!a{derivation} दिली आहे:
\begin{prooftree}
\iftag{prvBox}{\iftag{prvDiamond}{% both [] and <> primitive
\Axiom$\Box !A \fCenter \Box!A$
\RightLabel{5$\Diamond$}
\UnaryInf$\Diamond\Box!A \fCenter \Box!A$

\Axiom$!A \fCenter !A$
\RightLabel{T$\Box$}
\UnaryInf$\Box!A \fCenter !A$
\RightLabel{\Cut}
\BinaryInf$\Diamond\Box!A \fCenter !A$
\RightLabel{\RightR{\lif}}
\UnaryInf$\fCenter \Diamond\Box!A \lif !A$
}{% only [] primitive
\Axiom$\Box !A \fCenter \Box!A$
\RightLabel{\RightR{\lnot}}
\UnaryInf$\fCenter \Box!A, \lnot\Box!A$
\RightLabel{5$\Box$}
\UnaryInf$\fCenter \Box!A, \Box\lnot\Box!A$
\RightLabel{\RightR{\Exchange}}
\UnaryInf$\fCenter \Box\lnot\Box!A, \Box!A$

\Axiom$!A \fCenter !A$
\RightLabel{T$\Box$}
\UnaryInf$\Box!A \fCenter !A$
\RightLabel{\Cut}
\BinaryInf$ \fCenter \Box\lnot\Box!A, !A$
\RightLabel{\LeftR{\lnot}}
\UnaryInf$\lnot\Box\lnot\Box!A \fCenter !A$
\RightLabel{\RightR{\lif}}
\UnaryInf$\fCenter \lnot\Box\lnot\Box!A \lif !A$
}}{% only <> primitive
\Axiom$!A \fCenter !A$
\RightLabel{T$\Diamond$}
\UnaryInf$!A \fCenter \Diamond!A$

\Axiom$\Diamond !A \fCenter \Diamond!A$
\RightLabel{\LeftR{\lnot}}
\UnaryInf$\lnot\Diamond!A , \Diamond!A \fCenter$
\RightLabel{5$\Diamond$}
\UnaryInf$\Diamond\lnot\Diamond!A , \Diamond!A\fCenter$
\RightLabel{\LeftR{\Exchange}}
\UnaryInf$\Diamond!A, \Diamond\lnot\Diamond!A\fCenter $

\RightLabel{\Cut}
\BinaryInf$\Diamond\lnot\Diamond!A, !A\fCenter $
\RightLabel{\RightR{\lnot}}
\UnaryInf$!A \fCenter \lnot\Diamond\lnot\Diamond!A$
\RightLabel{\RightR{\lif}}
\UnaryInf$\fCenter !A \lif \lnot\Diamond\lnot\Diamond!A$
}
\end{prooftree}
\end{ex}

\begin{prob}
पुढील निष्कर्ष दाखवणाऱ्या क्रमवर्ती !!{derivation}s द्या:
  \begin{enumerate}
  \item $\Log{KT5} \Proves \Ax{B}$;
  \item $\Log{KT5} \Proves \Ax{4}$;
  \item $\Log{KDB4} \Proves \Ax{T}$;
  \item $\Log{KB4} \Proves \Ax{5}$;
  \item $\Log{KB5} \Proves \Ax{4}$;
  \item $\Log{KT} \Proves \Ax{D}$.
  \end{enumerate}
\end{prob}

\end{document}
